\documentclass[12pt,a4paper]{article}

% ── Paquetes ─────────────────────────────────────────────────────────────────
\usepackage[utf8]{inputenc}
\usepackage[T1]{fontenc}
\usepackage[spanish]{babel}
\usepackage{geometry}
\usepackage{amsmath}
\usepackage{booktabs}
\usepackage{array}
\usepackage{longtable}
\usepackage{xcolor}
\usepackage{listings}
\usepackage{hyperref}
\usepackage{enumitem}
\usepackage{titlesec}
\usepackage{fancyhdr}
\usepackage{graphicx}
\usepackage{microtype}
\usepackage{parskip}

\geometry{margin=2.5cm}

% ── Colores ───────────────────────────────────────────────────────────────────
\definecolor{codeblue}{RGB}{0,90,160}
\definecolor{codegray}{RGB}{100,100,100}
\definecolor{codered}{RGB}{160,30,30}
\definecolor{codegreen}{RGB}{0,120,60}
\definecolor{backcolor}{RGB}{248,248,248}

% ── Estilo de código ──────────────────────────────────────────────────────────
\lstdefinestyle{cstyle}{
  backgroundcolor=\color{backcolor},
  basicstyle=\ttfamily\small,
  keywordstyle=\color{codeblue}\bfseries,
  commentstyle=\color{codegray}\itshape,
  stringstyle=\color{codered},
  breaklines=true,
  numbers=left,
  numberstyle=\tiny\color{codegray},
  frame=single,
  framerule=0.4pt,
  rulecolor=\color{codegray},
  tabsize=4,
  language=C
}
\lstset{style=cstyle}

% ── Encabezado y pie ─────────────────────────────────────────────────────────
\pagestyle{fancy}
\fancyhf{}
\rhead{\small Sistemas Distribuidos y Paralelos}
\lhead{\small Práctica 3 — MPI}
\cfoot{\thepage}
\renewcommand{\headrulewidth}{0.4pt}

% ── Formato de secciones ──────────────────────────────────────────────────────
\titleformat{\section}{\large\bfseries\color{codeblue}}{Ejercicio \thesection.}{0.5em}{}[\titlerule]
\titleformat{\subsection}{\normalsize\bfseries}{}{0em}{}

% ── Tabla auxiliar ────────────────────────────────────────────────────────────
\newcolumntype{L}[1]{>{\raggedright\arraybackslash}p{#1}}

% ─────────────────────────────────────────────────────────────────────────────
\begin{document}

% ── Portada ──────────────────────────────────────────────────────────────────
\begin{titlepage}
  \centering
  \vspace*{2cm}
  {\Huge\bfseries Práctica 3\par}
  \vspace{0.5cm}
  {\LARGE Programación Paralela con MPI\par}
  \vspace{1.5cm}
  {\large Sistemas Distribuidos y Paralelos\par}
  \vspace{3cm}
  \begin{abstract}
    \noindent
    Este informe documenta el diseño e implementación de soluciones paralelas usando
    MPI (Message Passing Interface) en C para los algoritmos secuenciales de la
    Práctica 1. Por cada ejercicio se aplica la metodología de diseño de Foster
    (Descomposición, Comunicación, Aglomeración, Mapeo), se analizan las ventajas
    y desventajas de las herramientas utilizadas, y se estudian alternativas en
    otros lenguajes de programación tanto compilados como interpretados.
  \end{abstract}
  \vfill
  {\large \today\par}
\end{titlepage}

% ── Índice ───────────────────────────────────────────────────────────────────
\tableofcontents
\newpage

% =============================================================================
\section*{Marco Teórico: Metodología de Diseño Foster}
\addcontentsline{toc}{section}{Marco Teórico: Metodología de Diseño Foster}
% =============================================================================

Todos los ejercicios de esta práctica siguen la \textbf{metodología de diseño Foster}
para la paralelización sistemática de algoritmos:

\begin{enumerate}[leftmargin=*]
  \item \textbf{Descomposición}: identificar las unidades de trabajo independientes.
    Puede ser \textit{de datos} (se dividen los datos entre procesos) o
    \textit{exploratoria} (se divide el espacio de búsqueda).
  \item \textbf{Comunicación}: determinar qué datos deben compartirse entre procesos
    y con qué topología (broadcast, scatter/gather, all-to-all, punto a punto).
  \item \textbf{Aglomeración}: agrupar tareas finas en bloques más gruesos para
    amortizar el overhead fijo de latencia de cada mensaje MPI.
  \item \textbf{Mapeo}: asignar bloques de trabajo a procesos. Puede ser
    \textit{estático} (asignación fija antes de la ejecución) o
    \textit{dinámico} (asignación en tiempo de ejecución según disponibilidad).
\end{enumerate}

El patrón dominante en los ejercicios 1, 2, 3 y 5 es \textbf{SPMD} (\textit{Single
Program Multiple Data}): todos los procesos ejecutan el mismo código y se bifurcan
únicamente con \texttt{if (rank == 0)} para separar el rol coordinador del rol
computador. El ejercicio 4 usa el patrón \textbf{Master-Worker} con distribución
dinámica de tareas (\textit{Bag of Tasks}).

\newpage

% =============================================================================
\section{Multiplicación de Matrices con MPI}
% =============================================================================

\subsection{Diseño Foster}

\begin{tabular}{L{3cm} L{11cm}}
\toprule
\textbf{Fase} & \textbf{Descripción} \\
\midrule
Descomposición & De datos por filas de $A$. La fila $i$ de $C = A \times B$
  solo depende de la fila $i$ de $A$ y de toda $B$ $\Rightarrow$ las $N$ filas de $C$
  son completamente independientes entre sí. \\[4pt]
Comunicación & \texttt{MPI\_Bcast} de $B$ completa ($O(N^2)$ datos),
  \texttt{MPI\_Scatter} de $N/P$ filas de $A$, \texttt{MPI\_Gather} de $N/P$
  filas de $C$. \\[4pt]
Aglomeración & Bloque de $N/P$ filas contiguas por proceso. Granularidad
  gruesa que minimiza el overhead de comunicación relativo al cómputo. \\[4pt]
Mapeo & Estático. Cada proceso realiza exactamente $(N/P) \cdot N^2$
  operaciones de punto flotante $\Rightarrow$ carga perfectamente uniforme. \\
\bottomrule
\end{tabular}

\vspace{0.5cm}

\subsection{Variantes implementadas}

\textbf{a) mm\_naive optimizado (p3Ej1a.c):} El kernel local usa el orden de
bucles \texttt{ikj} en lugar del clásico \texttt{ijk}. Con orden \texttt{ijk},
el acceso a $B$ es por columnas (saltos de $N$ posiciones en memoria), generando
\textit{cache misses} continuos. Con \texttt{ikj}, el factor $a_{ik}$ queda fijo
mientras se recorre la fila $j$ de $B$ de forma secuencial, aprovechando las
líneas de caché.

\textbf{b) Autoproducto $C = A \cdot A$ (p3Ej1b.c):} La variante donde ambas
operandas son la misma matriz. Se realiza un \texttt{MPI\_Bcast} de $A$ completa
(para usarla como ``segunda operanda'') más un \texttt{MPI\_Scatter} de $A$ (para
las filas locales). La validación es \textit{inline}: rank 0 computa $A \times A$
secuencialmente con el mismo algoritmo \texttt{ikj} y compara con el resultado
paralelo, sin necesitar un archivo de referencia externo.

\textbf{c) Multiplicación por bloques (p3Ej1c.c):} La comunicación MPI es idéntica
a la variante naive. La diferencia está en el kernel local: en lugar de procesar filas
completas, se trabaja con sub-bloques de tamaño $\text{BS} \times \text{BS}$. Si
$\text{BS}$ es apropiado para la jerarquía de caché, el sub-bloque permanece en L1/L2
durante toda su utilización, reduciendo los \textit{cache misses} de forma
significativa (principio de localidad temporal y espacial). La variante compilada con
\texttt{-DUSE\_BLAS} aplica adicionalmente una transposición local del bloque de $B$
para garantizar acceso secuencial a memoria.

\subsection{Ventajas y desventajas de MPI en C}

\begin{longtable}{L{4cm} L{10cm}}
\toprule
\textbf{Ventajas} & \textbf{Detalle} \\
\midrule
\endhead
Control total de memoria & El programador decide exactamente qué datos residen en
  cada proceso, evitando copias innecesarias. Crítico para matrices grandes. \\[4pt]
Portabilidad & MPI funciona en clusters heterogéneos, supercomputadoras con
  InfiniBand, y máquinas locales con Ethernet. Una misma compilación puede
  ejecutarse en cualquier plataforma que tenga una implementación MPI
  (OpenMPI, MPICH, Intel MPI). \\[4pt]
Rendimiento predecible & La latencia y el \textit{throughput} de cada comunicación
  son explícitos y medibles con \texttt{MPI\_Wtime}. No hay abstracciones ocultas. \\[4pt]
Integración con BLAS/LAPACK & El código C puede llamar directamente a librerías
  altamente optimizadas para el cómputo local (dgemm, etc.) sin capas
  intermedias de conversión. \\
\midrule
\textbf{Desventajas} & \\
\midrule
Verbosidad & Inicializar comunicadores, calcular desplazamientos, gestionar
  manualmente la memoria local de cada proceso y manejar el caso $N \% P \neq 0$
  requiere decenas de líneas de código de ``infraestructura''. \\[4pt]
Debugging difícil & Un deadlock o una condición de carrera en comunicaciones
  puede colgar silenciosamente el programa. Las herramientas (Valgrind/MPI,
  TotalView) son complejas de usar. \\[4pt]
Sin tolerancia a fallos & Si un proceso muere, el job completo falla. MPI no
  tiene recuperación nativa ante fallos de nodo. \\[4pt]
Manejo manual del balance & Si los bloques no son uniformes (N no divisible por P),
  hay que implementar lógica adicional con \texttt{MPI\_Scatterv}. \\
\bottomrule
\end{longtable}

\subsection{Alternativas en otros lenguajes}

\begin{longtable}{L{2.5cm} L{3.5cm} L{7.5cm}}
\toprule
\textbf{Lenguaje} & \textbf{Herramienta} & \textbf{Análisis} \\
\midrule
\endhead
Python (interpretado) & \texttt{mpi4py} + NumPy & MPI con sintaxis Pythónica. Las comunicaciones de
  arrays NumPy evitan la serialización y se traducen directamente a operaciones
  MPI en C subyacente. El overhead del intérprete es insignificante si el
  cómputo local usa NumPy (\texttt{np.dot}) o SciPy (que llaman BLAS en C).
  Ideal para prototipado rápido. Para producción a gran escala, el overhead del
  intérprete en Python puro puede ser un factor 10--100x más lento que C. \\[4pt]
Julia (compilado JIT) & \texttt{MPI.jl} & Julia compila JIT a código nativo
  comparable a C. El paquete \texttt{MPI.jl} provee bindings completos a la
  API de MPI. Es la alternativa más cercana a C en rendimiento con la
  ergonomía de un lenguaje de alto nivel: manejo automático de tipos, arrays
  multidimensionales nativos, y broadcasting. \\[4pt]
Rust (compilado) & \texttt{rsmpi} & Rust garantiza ausencia de carreras de
  datos en tiempo de compilación mediante su sistema de \textit{ownership}.
  La librería \texttt{rsmpi} expone MPI con seguridad de tipos. Para código
  nuevo donde la correctitud es crítica, Rust elimina una clase entera de bugs.
  La curva de aprendizaje es empinada. \\[4pt]
Java (compilado JVM) & MPJ Express, FastMPJ & Java elimina la gestión manual de
  memoria pero el JVM agrega overhead de GC y boxing de tipos primitivos. Para
  matrices de \texttt{double[]}, el rendimiento es aceptable pero nunca comparable
  a C. Útil en entornos donde el ecosistema Java ya existe. \\[4pt]
Go (compilado) & \texttt{gompi} & Go ofrece goroutines para paralelismo de
  memoria compartida y \texttt{gompi} para paso de mensajes. Compila a nativo.
  Para computación científica pura su ecosistema de álgebra lineal es más
  limitado que el de C o Julia. \\
\bottomrule
\end{longtable}

\textbf{Recomendación para matrices:} Para investigación/prototipado: Python con
\texttt{mpi4py} + NumPy. Para producción en HPC: C/MPI + BLAS o Julia + \texttt{MPI.jl}.

\newpage

% =============================================================================
\section{Estadísticas de Matrices: Mínimo, Máximo y Promedio}
% =============================================================================

\subsection{Diseño Foster}

\begin{tabular}{L{3cm} L{11cm}}
\toprule
\textbf{Fase} & \textbf{Descripción} \\
\midrule
Descomposición & De datos. La matriz se trata como un arreglo lineal de $N^2$
  elementos. Cada elemento contribuye de forma independiente al mínimo, máximo
  y suma $\Rightarrow$ $N^2$ tareas atómicas completamente independientes. \\[4pt]
Comunicación & \texttt{MPI\_Scatter} de $N^2/P$ elementos, luego
  \texttt{MPI\_Reduce} con \texttt{MPI\_MIN}, \texttt{MPI\_MAX} y \texttt{MPI\_SUM}.
  La comunicación de resultados parciales es $O(P)$, insignificante respecto
  al scatter de datos. \\[4pt]
Aglomeración & Bloque de $N^2/P$ elementos contiguos por proceso. Granularidad
  gruesa: una sola llamada de scatter distribuye todo el trabajo. \\[4pt]
Mapeo & Estático. Carga perfectamente uniforme: todos los procesos calculan
  exactamente $N^2/P$ comparaciones. \\
\bottomrule
\end{tabular}

\vspace{0.5cm}

\subsection{Patrón clave: MPI\_Reduce}

\texttt{MPI\_Reduce} es la operación colectiva que combina valores escalares de todos
los procesos usando una operación predefinida (\texttt{MPI\_MIN}, \texttt{MPI\_MAX},
\texttt{MPI\_SUM}, \texttt{MPI\_PROD}, etc.) y deposita el resultado en el proceso
raíz. Internamente utiliza un árbol binario de comunicaciones en $O(\log P)$ pasos,
siendo significativamente más eficiente que que rank 0 recibiera uno a uno de cada
proceso ($O(P)$ pasos seriales).

\subsection{Sobre el uso de \texttt{srand(42)}}

La inicialización de la matriz se realiza con \texttt{srand(42)}, una semilla fija
para el generador de números pseudo-aleatorios. Esto es intencional y tiene tres
propósitos:

\begin{itemize}
  \item \textbf{Reproducibilidad}: cada ejecución del programa genera exactamente los
    mismos datos. Un resultado diferente entre corridas indica un bug, no datos distintos.
  \item \textbf{Validación automática}: rank 0 puede re-recorrer el mismo arreglo
    secuencialmente (sin necesidad de guardarlo en disco) y comparar contra el resultado
    paralelo, ya que la secuencia de \texttt{rand()} es determinista dado el mismo seed.
  \item \textbf{Experimentos comparables}: al medir tiempo con $P = 1, 2, 4, 8, \ldots$
    procesos, todos los experimentos operan sobre exactamente los mismos datos,
    haciendo que la comparación de speedups sea justa.
\end{itemize}

En producción real se usaría \texttt{srand(time(NULL))} o un generador
criptográficamente seguro, pero para ejercicios de validación académica la
semilla fija es la práctica estándar.

\subsection{Ventajas y desventajas}

\begin{longtable}{L{4cm} L{10cm}}
\toprule
\textbf{Ventajas} & \textbf{Detalle} \\
\midrule
\endhead
Mínima comunicación & Solo se transfieren $P$ valores escalares en el Reduce
  $\Rightarrow$ eficiencia de comunicación muy alta. \\[4pt]
Escalabilidad ideal & El problema es \textit{embarazosamente paralelo}: no hay
  dependencias entre elementos. El speedup teórico se aproxima a $P$. \\[4pt]
Operaciones MPI nativas & \texttt{MPI\_MIN}, \texttt{MPI\_MAX}, \texttt{MPI\_SUM}
  son primitivas que los MPI modernos implementan con hardware especializado
  (reducción en la red InfiniBand, por ejemplo). \\
\midrule
\textbf{Desventajas} & \\
\midrule
Scatter de datos & El scatter inicial transfiere $O(N^2)$ datos. Para matrices muy
  grandes, esta distribución inicial puede dominar el tiempo total si el cómputo
  por elemento es trivial (solo comparaciones). \\[4pt]
Alternativa más simple & Para este problema específico, un enfoque de memoria
  compartida (OpenMP) sería más eficiente en una sola máquina, ya que evita el
  scatter inicial: la matriz ya está en memoria accesible por todos los hilos. \\
\bottomrule
\end{longtable}

\subsection{Alternativas en otros lenguajes}

\begin{longtable}{L{2.5cm} L{3.5cm} L{7.5cm}}
\toprule
\textbf{Lenguaje} & \textbf{Herramienta} & \textbf{Análisis} \\
\midrule
\endhead
Python & \texttt{mpi4py} + NumPy & \texttt{np.min()}, \texttt{np.max()},
  \texttt{np.mean()} sobre el chunk local son $O(N^2/P)$ con código C subyacente.
  El Reduce de \texttt{mpi4py} funciona directamente sobre escalares Python.
  Código extremadamente conciso (5--10 líneas vs. 90 en C). \\[4pt]
R (interpretado) & \texttt{Rmpi} o \texttt{pbdMPI} & R es un lenguaje diseñado
  para estadística. Sus funciones \texttt{min()}, \texttt{max()}, \texttt{mean()}
  son vectorizadas en C internamente. \texttt{pbdMPI} ofrece una API de MPI
  simplificada. Muy productivo para análisis de datos; no competitivo para
  HPC puro. \\[4pt]
Julia & \texttt{MPI.jl} & Sintaxis idéntica al problema: \texttt{minimum(local)},
  \texttt{MPI.Reduce} con \texttt{MPI.MIN}. Rendimiento equivalente a C. \\[4pt]
Fortran (compilado) & MPI nativo & Fortran tiene soporte MPI de primera clase y
  sigue siendo el estándar en computación científica numérica. Sus arrays
  multidimensionales son nativos (no punteros) y sus compiladores (gfortran,
  ifort) generan código altamente optimizado para operaciones numéricas. \\
\bottomrule
\end{longtable}

\newpage

% =============================================================================
\section{Transposición de Matriz $N \times N$}
% =============================================================================

\subsection{Diseño Foster}

\begin{tabular}{L{3cm} L{11cm}}
\toprule
\textbf{Fase} & \textbf{Descripción} \\
\midrule
Descomposición & De datos. Cada elemento $(i,j)$ de $A$ produce $T[j][i]$ $\Rightarrow$
  $N^2$ tareas independientes. Pero la dificultad está en la comunicación:
  el proceso que tiene la fila $i$ de $A$ debe entregarla a quien construye
  la columna $i$ de $T$ (que está en otro proceso). \\[4pt]
Comunicación & Redistribución total: \textbf{all-to-all}. Cada proceso envía datos a
  todos los demás. Se usa \texttt{MPI\_Alltoall} con un paso previo de
  empaquetado de sub-bloques. \\[4pt]
Aglomeración & Cada proceso maneja $N/P$ filas de $A$ y produce $N/P$ filas de $T$.
  Los sub-bloques de $N/P \times N/P$ elementos se intercambian via Alltoall. \\[4pt]
Mapeo & Estático. Carga uniforme: cada proceso envía y recibe exactamente
  $N^2/P$ elementos. \\
\bottomrule
\end{tabular}

\vspace{0.5cm}

\subsection{Algoritmo de empaquetado/desempaquetado}

El paso no trivial del algoritmo es el empaquetado del \texttt{send\_buf} antes del
Alltoall y el desempaquetado del \texttt{recv\_buf} después. Si el proceso $p$ tiene
las filas $[p \cdot ln, (p+1) \cdot ln - 1]$ de $A$ (donde $ln = N/P$), para
enviárselas al proceso $q$ necesita extraer el sub-bloque de $ln \times ln$ elementos
correspondiente a las columnas $[q \cdot ln, (q+1) \cdot ln - 1]$:

\[
  \texttt{send\_buf}[q \cdot ln^2 + i \cdot ln + k] = A[i][q \cdot ln + k]
  \quad \forall\, i,k \in [0, ln)
\]

Tras el Alltoall, el proceso $p$ recibe en \texttt{recv\_buf} los sub-bloques
provenientes de cada proceso $q$, que corresponden a $A[q \cdot ln + i][p \cdot ln + k]$
$= T[p \cdot ln + k][q \cdot ln + i]$, es decir, elementos de las filas locales de $T$.

\subsection{Patrón clave: MPI\_Alltoall}

\texttt{MPI\_Alltoall} es la operación colectiva más costosa en términos de
comunicación: cada uno de los $P$ procesos envía un chunk diferente a cada otro
proceso, totalizando $P^2$ mensajes implícitos. En términos de datos movidos,
$O(N^2)$ elementos en total, distribuidos de forma simétrica. Es la operación
apropiada cuando hay \textbf{dependencias cruzadas totales} entre la distribución
inicial y la distribución final de los datos.

\subsection{Ventajas y desventajas}

\begin{longtable}{L{4cm} L{10cm}}
\toprule
\textbf{Ventajas} & \textbf{Detalle} \\
\midrule
\endhead
Expresividad & \texttt{MPI\_Alltoall} encapsula en una sola llamada lo que
  sería un doble bucle de \texttt{MPI\_Send}/\texttt{MPI\_Recv} entre todos
  los pares, permitiendo al runtime optimizar el patrón de comunicación. \\[4pt]
Escalabilidad & El runtime puede usar algoritmos de Alltoall optimizados para
  la topología de la red (torus, árbol gordo, InfiniBand), imposibles de
  replicar manualmente. \\
\midrule
\textbf{Desventajas} & \\
\midrule
Costo de comunicación & Es la operación colectiva de mayor volumen. Con $P$
  procesos y matrices grandes, la red puede convertirse en el cuello de botella. \\[4pt]
Empaquetado manual & La API de \texttt{MPI\_Alltoall} exige que los datos
  estén contiguos en el buffer, por lo que el programador debe implementar
  el empaquetado y desempaquetado de sub-bloques, lo que agrega complejidad
  y posibilidad de errores de índices. \texttt{MPI\_Alltoallv} o tipos de datos
  derivados (\texttt{MPI\_Type\_vector}) pueden evitar la copia manual pero
  aumentan la complejidad del código. \\[4pt]
Barrera implícita & \texttt{MPI\_Alltoall} es síncrona: todos los procesos deben
  llegar a la llamada antes de que alguno pueda continuar. \\
\bottomrule
\end{longtable}

\subsection{Alternativas en otros lenguajes}

\begin{longtable}{L{2.5cm} L{3.5cm} L{7.5cm}}
\toprule
\textbf{Lenguaje} & \textbf{Herramienta} & \textbf{Análisis} \\
\midrule
\endhead
Python & \texttt{mpi4py} & \texttt{comm.Alltoall(sendbuf, recvbuf)} funciona directamente
  sobre arrays NumPy. El empaquetado puede hacerse con operaciones de slicing
  de NumPy (\texttt{arr[i::stride, :]}) que son significativamente más
  expresivas que el doble bucle en C. \\[4pt]
Julia & \texttt{MPI.jl} & Julia tiene indexing multidimensional nativo que hace que
  el empaquetado sea más legible. El rendimiento es equivalente a C. \\[4pt]
Chapel (compilado) & Nativo & Chapel es un lenguaje académico diseñado específicamente
  para HPC. Sus distribuciones de arreglos (\texttt{Block}, \texttt{Cyclic})
  abstraen completamente la distribución de datos: \texttt{var T = transpose(A)}
  puede ejecutarse en paralelo sin que el programador gestione comunicaciones. \\[4pt]
MATLAB/Octave & Parallel Computing Toolbox & \texttt{A'} traspone y el toolbox
  puede distribuirlo con \texttt{codistributed}. Solo viable para prototipado
  académico, no para producción. \\
\bottomrule
\end{longtable}

\newpage

% =============================================================================
\section{N-Reinas Paralelo (Master-Worker / Bag of Tasks)}
% =============================================================================

\subsection{Diseño Foster}

\begin{tabular}{L{3cm} L{11cm}}
\toprule
\textbf{Fase} & \textbf{Descripción} \\
\midrule
Descomposición & \textbf{Exploratoria}. El espacio de búsqueda es un árbol de
  backtracking. Se cortan ramas en los niveles 1--2 del árbol para definir
  tareas independientes (subtárboles). Dos tipos de tarea:
  (1) primera reina en esquina, segunda en columna $b_1$ (Backtrack1);
  (2) primera reina en posición interior $b_1$ (Backtrack2). \\[4pt]
Comunicación & Dinámica. Master envía una tarea $\{tipo, b_1, b_2\}$ por vez a
  cada worker mediante \texttt{MPI\_Send}, worker devuelve contadores
  \{COUNT8, COUNT4, COUNT2\} con \texttt{MPI\_Send}. Señal de fin con
  \texttt{TAG\_DONE}. \\[4pt]
Aglomeración & Un subtárbol completo por tarea $\Rightarrow$ granularidad
  media-alta. Evita overhead excesivo de comunicación si las tareas fueran
  muy finas (por ejemplo, una sola colocación de reina). \\[4pt]
Mapeo & \textbf{Dinámico centralizado} (Bag of Tasks). El master asigna tareas
  a demanda: cuando un worker termina una, recibe la siguiente disponible. \\
\bottomrule
\end{tabular}

\vspace{0.5cm}

\subsection{a) Estrategia de descomposición más adecuada}

Para el problema N-Reinas la estrategia más adecuada es la \textbf{descomposición
exploratoria}, que consiste en dividir el árbol de backtracking en subtárboles
independientes.

La razón por la que una descomposición de datos convencional no aplica es que no
existe un arreglo de entrada que distribuir: el ``trabajo'' está completamente
implícito en las posibles combinaciones del tablero. Cada subtárbol puede explorarse
de forma completamente independiente (no hay estado compartido entre ramas), lo que
los convierte en tareas MPI naturales.

La granularidad óptima se obtiene fijando las primeras 1--2 reinas del tablero:
\begin{itemize}
  \item Si se fijan muy pocas reinas (granularidad fina), el overhead de comunicación
    por tarea es proporcionalmente alto.
  \item Si se fijan muchas reinas (granularidad gruesa), el número de tareas se
    reduce y puede haber desbalance si algunas ramas grandes caen en el mismo proceso.
  \item Fijar las primeras 2 reinas genera entre 15 y 30 tareas para $N = 13$--$15$,
    un número razonable para balancear con unos pocos workers.
\end{itemize}

\subsection{b) ¿Qué es una tarea? ¿Son uniformes? Problema del mapeo estático}

Una \textbf{tarea} en este contexto es la exploración completa de un subtárbol del
árbol de backtracking, definida por la posición fija de las primeras 1--2 reinas del
tablero:

\begin{itemize}
  \item \textbf{Tipo 1 (Backtrack1)}: la primera reina está en la esquina ($\text{col} = 0$),
    la segunda en la columna $b_1 \in [2, N-2]$. Genera $N-4$ tareas.
  \item \textbf{Tipo 2 (Backtrack2)}: la primera reina está en la posición interior
    $b_1$, explorando la simetría por reflexión. Genera $\lfloor N/2 \rfloor - 1$ tareas.
\end{itemize}

Las tareas \textbf{no son uniformes}. El tiempo de ejecución de cada tarea depende
de cuántas ramas del subtárbol se podan y de cuántas soluciones contiene. Algunas
posiciones iniciales generan rápidamente contradicciones (poda agresiva, tarea barata)
mientras que otras habilitan miles de soluciones intermedias (tarea cara). La varianza
entre tareas puede ser de varios órdenes de magnitud.

\textbf{Problema del mapeo estático}: si se divide la lista de tareas en $P$ bloques
consecutivos asignados de forma fija a cada proceso, el tiempo total queda determinado
por el proceso que recibió las tareas más costosas. Un proceso puede quedar con múltiples
tareas de tiempo exponencial mientras otro termina todas las suyas en milisegundos y
queda inactivo. Esto puede hacer que el speedup real sea de $1.2\times$ en lugar del
$P\times$ teórico.

El \textbf{mapeo dinámico} (Bag of Tasks) resuelve este problema: el proceso más rápido
toma automáticamente más tareas de la bolsa, garantizando que ningún worker quede ocioso
mientras haya trabajo disponible.

\subsection{c) Implementación Master-Worker}

La implementación utiliza comunicación punto a punto con tres tipos de mensajes
identificados por tags:

\begin{itemize}
  \item \texttt{TAG\_TASK} (1): master $\to$ worker, contiene el descriptor de tarea
    $\{tipo, b_1, b_2\}$ como arreglo de 3 enteros.
  \item \texttt{TAG\_RESULT} (2): worker $\to$ master, contiene los contadores
    $\{COUNT8, COUNT4, COUNT2\}$ como arreglo de 3 \texttt{long int}.
  \item \texttt{TAG\_DONE} (3): master $\to$ workers, señal de terminación
    (tarea vacía con \texttt{tipo = -1}).
\end{itemize}

El master usa \texttt{MPI\_Recv(MPI\_ANY\_SOURCE)} para recibir del primer worker
que termine, garantizando que no haya espera innecesaria. Los resultados parciales
se acumulan con una suma simple ya que \textit{COUNT8, COUNT4, COUNT2} son
independientes entre subtárboles.

\subsection{Ventajas y desventajas}

\begin{longtable}{L{4cm} L{10cm}}
\toprule
\textbf{Ventajas} & \textbf{Detalle} \\
\midrule
\endhead
Balance de carga dinámico & El patrón Bag of Tasks garantiza que ningún proceso
  quede ocioso mientras haya tareas disponibles, siendo óptimo para cargas
  no uniformes. \\[4pt]
Simplicidad del modelo & Los workers son completamente \textit{stateless}: no
  necesitan conocer el estado global ni coordinarse entre sí. Cada worker
  recibe una tarea, la ejecuta y devuelve el resultado. \\[4pt]
Escalabilidad hasta cierto punto & Escala bien mientras la velocidad de
  distribución del master supere la velocidad de consumo de los workers. \\
\midrule
\textbf{Desventajas} & \\
\midrule
Cuello de botella en el master & El master no computa ninguna tarea: dedica
  todo su tiempo a coordinar. Con miles de workers, el master puede convertirse
  en el factor limitante. \\[4pt]
Latencia de granularidad fina & Si las tareas son muy cortas (microsegundos),
  el overhead de un round-trip de mensajes MPI (Send + Recv) puede superar el
  tiempo de cómputo de la tarea misma. \\[4pt]
Punto único de fallo & Si el master falla, todo el trabajo se pierde. \\
\bottomrule
\end{longtable}

\subsection{Alternativas al patrón Master-Worker}

Además del patrón Master-Worker, existen otras estrategias de implementación
paralela aplicables a problemas de búsqueda y en general:

\begin{enumerate}[leftmargin=*]
  \item \textbf{SPMD con mapeo cíclico}: cada proceso toma las tareas con índices
    $\{rank, rank+P, rank+2P, \ldots\}$. No requiere comunicación durante la
    ejecución. El desbalance es menor que el estático contiguo (las tareas pesadas
    y livianas se intercalan), pero no es tan perfecto como el dinámico.

  \item \textbf{Work Stealing descentralizado}: cada proceso tiene su propia cola.
    Cuando un proceso termina su cola, ``roba'' tareas de la cola de otro proceso.
    No hay master centralizado $\Rightarrow$ escala a miles de procesos. Difícil
    de implementar en MPI puro; más natural con MPI-3 RMA (acceso remoto a memoria).

  \item \textbf{Master-Worker jerárquico}: múltiples masters, cada uno coordinando
    un subconjunto de workers. Reduce el cuello de botella del master único. Útil
    cuando $P > 100$.

  \item \textbf{Pipeline con workers especializados}: si las tareas tienen etapas
    diferenciadas (generación, exploración, verificación), cada tipo de worker
    realiza una etapa y pasa resultados al siguiente tipo. No aplica directamente
    a N-Reinas pero sí a problemas con fases definidas.
\end{enumerate}

\subsection{Alternativas en otros lenguajes}

\begin{longtable}{L{2.5cm} L{3.5cm} L{7.5cm}}
\toprule
\textbf{Lenguaje} & \textbf{Herramienta} & \textbf{Análisis} \\
\midrule
\endhead
Python & \texttt{multiprocessing.Pool} o \texttt{concurrent.futures} &
  Para un solo nodo, la API de alto nivel \texttt{Pool.map()} implementa
  Bag of Tasks automáticamente. Para múltiples nodos, \texttt{mpi4py} con
  el mismo patrón. La facilidad de desarrollo compensa el overhead para
  problemas de tamaño moderado. \\[4pt]
Erlang/Elixir (interpretado) & OTP (Actor Model) & El modelo de actores de
  Erlang es naturalmente un Bag of Tasks: los actores son procesos livianos
  que se comunican por mensajes. La tolerancia a fallos es nativa (un actor
  puede morir y reiniciarse sin afectar al sistema). Para búsqueda distribuida
  a gran escala, el modelo de actores es más robusto que MPI. \\[4pt]
Scala (compilado JVM) & Akka & Similar a Erlang pero sobre JVM. Akka Router
  implementa Round-Robin o Smallest-Mailbox dispatch, equivalentes al Bag of
  Tasks. El ecosistema Java/Scala es más amplio. \\[4pt]
Go & \texttt{goroutines} + channels & El modelo de concurrencia de Go con
  canales es conceptualmente similar al Bag of Tasks: un canal de tareas
  y múltiples goroutines trabajadoras. Limitado a un nodo; para clusters
  se necesitaría una capa de red adicional. \\
\bottomrule
\end{longtable}

\newpage

% =============================================================================
\section{Merge Sort Paralelo (Árbol Binario de Fusiones)}
% =============================================================================

\subsection{Diseño Foster}

\begin{tabular}{L{3cm} L{11cm}}
\toprule
\textbf{Fase} & \textbf{Descripción} \\
\midrule
Descomposición & De datos. El arreglo de $N$ enteros se divide en $P$ partes
  iguales de $N/P$ elementos $\Rightarrow$ $P$ tareas independientes de
  ordenamiento local. \\[4pt]
Comunicación & Árbol binario de reducción en $\log_2 P$ rondas:
  Ronda $r$: el proceso con $rank \mod 2^{r+1} = 0$ recibe de $rank + 2^r$
  y fusiona. El proceso que envía termina su participación.
  Comunicación estructurada y síncrona (\texttt{MPI\_Send/Recv}). \\[4pt]
Aglomeración & $N/P$ elementos por proceso para el ordenamiento local,
  luego fusiones de tamaño creciente ($2N/P$, $4N/P$, $\ldots$, $N$). \\[4pt]
Mapeo & Estático para el sort local (todos hacen $N/P$ elementos). En las
  fusiones la carga baja exponencialmente: los procesos de mayor rank terminan
  pronto y quedan inactivos. \\
\bottomrule
\end{tabular}

\vspace{0.5cm}

\subsection{a) Implementación y análisis}

El algoritmo de merge sort paralelo consta de dos fases:

\textbf{Fase 1 -- Ordenamiento local:} Cada proceso ordena su fragmento de $N/P$
elementos con \texttt{qsort}. Esta fase es completamente paralela: no hay comunicación
y la carga es perfectamente uniforme. Complejidad: $O\!\left(\frac{N}{P} \log \frac{N}{P}\right)$.

\textbf{Fase 2 -- Fusiones en árbol binario:} En cada ronda $r$ (de 0 a $\log_2 P - 1$),
la mitad de los procesos activos se convierten en ``enviadores'' y la otra mitad en
``fusionadores''. Un proceso envía cuando su rank satisface:
\[
  rank \bmod (2 \cdot 2^r) \neq 0 \quad \text{pero} \quad rank \bmod 2^r = 0
\]
El proceso que recibe fusiona dos arreglos ya ordenados en tiempo $O(n_1 + n_2)$.

\textbf{Análisis de la fase de fusiones:}
\begin{itemize}
  \item Ronda 0: $P/2$ pares fusionan, cada uno $2 \cdot N/P$ elementos. Cómputo total: $P/2 \cdot 2N/P = N$.
  \item Ronda 1: $P/4$ pares fusionan, cada uno $4 \cdot N/P$ elementos. Cómputo total: $N$.
  \item $\vdots$
  \item Ronda $\log_2 P - 1$: 1 proceso fusiona $N$ elementos. Cómputo: $N$.
\end{itemize}

Cada ronda tiene el mismo cómputo total $N$, pero recae sobre \textit{menos} procesos
en cada ronda. El cuello de botella es rank 0, que en la última ronda hace solo $N$
operaciones pero en secuencia (no paralelo). El tiempo de la fase de fusiones es
$O(N \log P / 1) = O(N \log P)$ en el proceso 0, que es el mismo orden que el sort
secuencial de referencia $O(N \log N)$ cuando $P$ es pequeño.

\textbf{¿Por qué punto a punto y no \texttt{MPI\_Gather}?}
\texttt{MPI\_Gather} acumula todos los fragmentos en rank 0, que tendría que fusionarlos
todos solo (costo $O(N)$ en la fusión final). La estructura en árbol distribuye el
trabajo de fusión entre múltiples procesos: en la ronda 0, la mitad de los procesos
fusionan en paralelo, reduciendo el trabajo total de fusión.

\subsection{Ventajas y desventajas}

\begin{longtable}{L{4cm} L{10cm}}
\toprule
\textbf{Ventajas} & \textbf{Detalle} \\
\midrule
\endhead
Sort local eficiente & \texttt{qsort} es altamente optimizado en la libc estándar.
  El sort local escala perfectamente con $P$. \\[4pt]
Sin comunicación en fase 1 & El sort local es completamente independiente
  entre procesos. \\[4pt]
Fusión en árbol distribuida & La carga de fusión se distribuye entre múltiples
  procesos en las primeras rondas, siendo más eficiente que un único proceso
  fusionando todo. \\
\midrule
\textbf{Desventajas} & \\
\midrule
Carga decreciente en fusiones & En las últimas rondas, la mayoría de los procesos
  están inactivos. El speedup real es menor que $P$. \\[4pt]
Memoria creciente en rank 0 & Rank 0 debe poder alocar el arreglo completo de
  $N$ elementos al final. En la implementación, todos los procesos alocan $N$
  enteros desde el inicio para simplicidad. \\[4pt]
Comunicación de tamaño creciente & Los mensajes de fusión crecen exponencialmente:
  $N/P, 2N/P, \ldots, N/2$. Los mensajes grandes pueden saturar la red. \\[4pt]
Solo funciona con $P$ potencia de 2 & El árbol binario simétrico requiere que $P$
  sea potencia de 2 para que la estructura de rondas sea correcta. \\
\bottomrule
\end{longtable}

\subsection{Alternativas en otros lenguajes}

\begin{longtable}{L{2.5cm} L{3.5cm} L{7.5cm}}
\toprule
\textbf{Lenguaje} & \textbf{Herramienta} & \textbf{Análisis} \\
\midrule
\endhead
Python & \texttt{mpi4py} + \texttt{heapq.merge} & El sort local puede hacerse con
  \texttt{sorted()} (implementado en C, muy eficiente). La fusión con
  \texttt{heapq.merge()} es $O(N \log P)$ por ser una fusión de $k$-vías en
  lugar de binaria, pero requiere solo una ronda de \texttt{MPI\_Gather} y
  una fusión final en rank 0. Para datos en memoria, más simple que el árbol. \\[4pt]
Java & \texttt{Arrays.sort()} + MPJ Express & Java 8+ usa TimSort, que es
  $O(N \log N)$ pero muy eficiente en datos parcialmente ordenados. La
  fusión en árbol se implementaría igual que en C. El boxing de primitivos
  (\texttt{int[]}) no aplica para arrays primitivos, por lo que el rendimiento
  es competitivo. \\[4pt]
Rust & \texttt{rayon} (local) + \texttt{rsmpi} (cluster) & En Rust la fase local
  puede usar \texttt{rayon::par\_sort()} que aplica quicksort paralelo en memoria
  compartida automáticamente, siendo más rápida que qsort en la fase 1. La
  fase de fusiones se implementa con \texttt{rsmpi}. Combina paralelismo de
  memoria compartida e intercomunicación de nodos. \\[4pt]
Spark/Hadoop (distribuido) & \texttt{RDD.sortBy()} & Para volúmenes de datos que
  no caben en RAM de un nodo, frameworks como Spark implementan ordenamiento
  distribuido con mucho menos código. El programador no gestiona MPI.
  Overhead significativo para pequeños $N$; ventajoso para $N > 10^9$. \\
\bottomrule
\end{longtable}

\newpage

% =============================================================================
\section*{Conclusiones}
\addcontentsline{toc}{section}{Conclusiones}
% =============================================================================

\subsection*{Receta general para paralelizar con MPI}

\begin{enumerate}[leftmargin=*, label=\textbf{Paso \arabic*.}]
  \item \textbf{¿El trabajo puede dividirse en partes independientes?}
    Sí $\to$ descomposición de datos (Ej1, Ej2, Ej3, Ej5).
    No $\to$ descomposición exploratoria (Ej4: árbol de backtracking).

  \item \textbf{¿La carga es uniforme?}
    Sí $\to$ mapeo estático, usar operaciones colectivas (Scatter/Gather/Reduce/Alltoall).
    No $\to$ mapeo dinámico, usar punto a punto (Master-Worker o Work Stealing).

  \item \textbf{¿Cuánta comunicación hay entre procesos?}
    \begin{itemize}
      \item Cada proceso solo necesita su parte $\to$ Scatter + cómputo + Gather.
      \item Un dato compartido por todos $\to$ agregar Bcast al inicio.
      \item Redistribución total $\to$ MPI\_Alltoall con empaquetado manual.
      \item Reducción jerárquica $\to$ árbol de MPI\_Send/Recv o MPI\_Reduce.
    \end{itemize}

  \item \textbf{Optimizar el kernel local antes que la comunicación.}
    La comunicación es una fracción menor del tiempo total solo si el cómputo
    local es eficiente. El cambio de orden de bucles (ikj vs ijk) o el blocking
    puede dar más ganancia que duplicar el número de procesos.

  \item \textbf{Medir con \texttt{MPI\_Wtime()} y validar contra resultado secuencial.}
    Siempre usar datos reproducibles (\texttt{srand} fijo) para poder comparar
    entre distintas configuraciones de $P$.
\end{enumerate}

\subsection*{Comparación de patrones MPI}

\begin{longtable}{L{3.5cm} L{2cm} L{2cm} L{2.5cm} L{3.5cm}}
\toprule
\textbf{Patrón} & \textbf{Escala} & \textbf{Balance} & \textbf{Comunicación} & \textbf{Cuándo usarlo} \\
\midrule
\endhead
SPMD / Data Parallel & Alta & Bueno & Colectivas & Datos uniformemente divisibles \\
Master-Worker & Media & Excelente & Punto a punto & Tareas no uniformes \\
Pipeline & Media & Limitado & Punto a punto & Fases diferenciadas \\
Árbol de reducción & Alta & Parcial & Punto a punto & Combinar resultados \\
All-to-All & Media & Uniforme & Alltoall & Redistribución total de datos \\
Work Stealing & Muy alta & Excelente & RMA / P2P & $P > 100$, tareas muy variables \\
\bottomrule
\end{longtable}

\subsection*{Selección de lenguaje}

\begin{itemize}
  \item \textbf{Producción HPC} (supercomputadoras, clusters InfiniBand): \textbf{C/MPI} o
    \textbf{Fortran/MPI}, potencialmente con BLAS/LAPACK para kernels numéricos.
  \item \textbf{Investigación y prototipado}: \textbf{Python + mpi4py + NumPy}. El código
    es 5--10x más conciso, el overhead es aceptable mientras NumPy realice el
    cómputo pesado en C internamente.
  \item \textbf{Rendimiento C + ergonomía moderna}: \textbf{Julia + MPI.jl}. Rendimiento
    equivalente a C con un lenguaje de nivel más alto y un ecosistema científico rico.
  \item \textbf{Correctitud crítica / sistemas embebidos}: \textbf{Rust + rsmpi}.
    Las garantías del compilador eliminan carreras de datos y errores de gestión
    de memoria en tiempo de compilación.
  \item \textbf{Big Data distribuido} ($N > 10^9$, datos en disco): \textbf{Apache Spark}
    o frameworks equivalentes; MPI no es el modelo adecuado para datos que no
    caben en RAM.
\end{itemize}

\end{document}
